\documentclass[11pt]{article}
\usepackage[utf8]{inputenc}
\usepackage[T1]{fontenc}
\usepackage{amsmath,amssymb}
\usepackage{graphicx}
\usepackage{booktabs}
\usepackage{hyperref}
\usepackage[margin=1in]{geometry}
\usepackage{natbib}
\usepackage{authblk}

\title{Noncanonical Inflammasome Assembly Requires Caspase-11 Catalytic Activity and Intra-Molecular Autoprocessing}

\author[1]{Author A}
\author[1]{Author B}
\author[1,2]{Author C}
\affil[1]{Department of Immunology, School of Medicine, University of Connecticut, Farmington, CT, USA}
\affil[2]{Institute for Systems Genomics, University of Connecticut, Farmington, CT, USA}

\date{}

\begin{document}
\maketitle

\begin{abstract}
Caspase-11 is the cytosolic sensor for bacterial lipopolysaccharide (LPS) and drives noncanonical inflammasome signaling and pyroptosis via gasdermin D (GSDMD) cleavage. While caspase-11 is known to form high molecular weight complexes upon LPS binding, whether it assembles into visible supramolecular organizing centers (SMOCs) in cells and whether catalytic activity is required for this assembly are unknown. Here, using fluorescently tagged caspase-11 constructs expressed in macrophages, we demonstrate that LPS transfection induces caspase-11 speck formation in 30--50\% of cells within 2--6 hours. Unexpectedly, both catalytic activity (C254A mutant) and intra-molecular autoprocessing (D285A mutant) are required for speck formation, as both mutants show severely reduced speck assembly ($<$5\% of cells). Biochemical characterization by size-exclusion chromatography confirms that catalytic-dead and autoprocessing-deficient caspase-11 fail to form high molecular weight oligomeric complexes. Both mutants also fail to cleave GSDMD and induce pyroptosis ($<$15\% LDH release versus 60--80\% for wild-type). These findings revise the prevailing model of inflammasome assembly: rather than being passively recruited to a pre-formed complex, caspase-11 enzymatic activity actively drives noncanonical inflammasome oligomerization.
\end{abstract}

\section{Introduction}

Inflammasomes are large oligomeric signaling complexes that activate inflammatory caspases, leading to cytokine maturation and an inflammatory form of cell death called pyroptosis \citep{broz2016inflammasomes}. Canonical inflammasomes assemble around sensor proteins such as NLRP3 or AIM2 that recruit the adaptor ASC, forming a single large perinuclear aggregate termed a speck or supramolecular organizing center (SMOC) \citep{lu2014unified}. Within this complex, proximity-induced dimerization activates caspase-1, which then cleaves pro-IL-1$\beta$, pro-IL-18, and gasdermin D (GSDMD) to execute pyroptosis \citep{shi2015cleavage}.

In contrast to the canonical pathway, the noncanonical inflammasome is driven by caspase-11 (caspase-4 and -5 in humans), which directly senses cytosolic LPS from Gram-negative bacteria \citep{shi2014inflammatory, kayagaki2011non}. LPS binding to the CARD domain of caspase-11 induces oligomerization and activation, leading to GSDMD cleavage and pyroptotic cell death \citep{shi2014inflammatory}. While biochemical studies have demonstrated that caspase-11 forms high molecular weight complexes upon LPS stimulation, direct visualization of caspase-11 specks in living cells has not been achieved. Furthermore, the prevailing model assumes that LPS-induced oligomerization occurs passively, with catalytic activation occurring subsequently via proximity-induced autoprocessing---analogous to the canonical inflammasome mechanism.

Here, we directly visualize LPS-induced caspase-11 speck formation in macrophages and make the unexpected discovery that catalytic activity and intra-molecular autoprocessing are required for SMOC assembly itself, not merely for downstream signaling. This fundamentally revises our understanding of noncanonical inflammasome assembly.

\section{Related Work}

The canonical inflammasome field has established the paradigm of sensor-driven oligomerization leading to caspase activation. \citet{lu2014unified} showed that ASC specks serve as SMOCs for caspase-1 activation, and subsequent studies demonstrated that these specks are single copy per cell and propagate in a prion-like manner \citep{franklin2014adaptor}. For the NLRP3 inflammasome, NEK7 was identified as an essential activating cofactor \citep{he2016nek7}.

Caspase-11 was identified as the sensor for cytosolic LPS by \citet{shi2014inflammatory}, who demonstrated direct binding between LPS and the caspase-11 CARD domain. \citet{kayagaki2011non} showed that caspase-11 is essential for LPS-induced septic shock independently of caspase-1. The downstream effector gasdermin D was identified by \citet{shi2015cleavage} and \citet{kayagaki2015caspase} as the executioner of pyroptosis.

While oligomerization of caspase-11 has been observed biochemically \citep{ross2018dimerization}, live-cell imaging of caspase-11 specks has not been reported. The question of whether catalytic activity is required for inflammasome assembly, as opposed to only downstream signaling, has been explored for caspase-1 in the canonical pathway \citep{broz2010redundant} but not for caspase-11 in the noncanonical context.

\section{Methods}

\subsection{Construct design}
Three caspase-11 constructs were generated with C-terminal GFP tags: (1) Casp11-WT-GFP (wild-type, catalytically active with intact autoprocessing); (2) Casp11-C254A-GFP (catalytic-dead mutant with the active-site cysteine mutated to alanine); and (3) Casp11-D285A-GFP (autoprocessing-deficient mutant with the interdomain cleavage site mutated). A GFP-only construct served as a negative control.

\subsection{Cell culture and LPS transfection}
Bone marrow-derived macrophages (BMDMs) were generated from C57BL/6 mice and transduced with lentiviral vectors encoding the Casp11 constructs. Cytosolic delivery of LPS was achieved by transfection with FuGENE HD, which bypasses the TLR4 pathway and directly engages the noncanonical inflammasome \citep{shi2014inflammatory}.

\subsection{Live-cell fluorescence microscopy}
Transduced macrophages were seeded on glass-bottom dishes and monitored by confocal fluorescence microscopy following LPS transfection. GFP speck formation was quantified as the percentage of GFP$^+$ cells containing at least one discrete fluorescent punctum at indicated time points.

\subsection{Pyroptosis assays}
GSDMD cleavage was assessed by Western blot using antibodies against the N-terminal fragment. Cell death was quantified by lactate dehydrogenase (LDH) release into the culture supernatant using a colorimetric assay. Propidium iodide (PI) uptake was assessed by flow cytometry and live imaging.

\subsection{Size-exclusion chromatography}
Cell lysates from LPS-stimulated or unstimulated macrophages expressing each construct were fractionated by size-exclusion chromatography on a Superose 6 column. Fractions were analyzed by Western blot to determine the oligomeric state of caspase-11.

\section{Results}

\subsection{LPS induces caspase-11 speck formation in macrophages}
Transfection of LPS into macrophages expressing Casp11-WT-GFP resulted in the formation of discrete fluorescent specks in 30--50\% of cells within 2--6 hours post-transfection (Table~\ref{tab:specks}). These specks were typically single-copy per cell and perinuclear, reminiscent of ASC specks in the canonical inflammasome. GFP-only expressing cells showed no speck formation, confirming that the specks are caspase-11-dependent.

\begin{table}[h]
\centering
\caption{LPS-induced speck formation by caspase-11 constructs.}
\label{tab:specks}
\begin{tabular}{lccc}
\toprule
Construct & Speck formation & \% Cells with specks & Timeframe \\
\midrule
Casp11-WT-GFP & Yes & 30--50\% & 2--6 h \\
Casp11-C254A-GFP & Severely reduced & $<$5\% & Up to 8 h \\
Casp11-D285A-GFP & Severely reduced & $<$5\% & Up to 8 h \\
GFP only & No & 0\% & -- \\
\bottomrule
\end{tabular}
\end{table}

\subsection{Catalytic activity and autoprocessing are required for speck formation}
The catalytic-dead mutant (C254A) and the autoprocessing-deficient mutant (D285A) both showed severely impaired speck formation ($<$5\% of cells, even at extended time points up to 8 hours; Table~\ref{tab:specks}). This was unexpected, as the prevailing model of inflammasome assembly posits that oligomerization occurs first and catalytic activity follows. Our data indicate that caspase-11 enzymatic activity is required upstream of or concurrently with SMOC assembly.

\subsection{Catalytic mutants fail to drive pyroptosis}
Consistent with the defect in speck formation, both catalytic-dead and autoprocessing-deficient mutants failed to cleave GSDMD, release LDH, or incorporate PI (Table~\ref{tab:pyroptosis}). Wild-type caspase-11 induced robust GSDMD cleavage, 60--80\% LDH release, and high PI uptake. Cells left untransfected with LPS showed minimal background death ($<$5\% LDH release).

\begin{table}[h]
\centering
\caption{Pyroptosis readouts for caspase-11 constructs after LPS transfection.}
\label{tab:pyroptosis}
\begin{tabular}{lccc}
\toprule
Construct & GSDMD cleavage & LDH release (\%) & PI uptake \\
\midrule
Casp11-WT & Yes (strong) & 60--80 & High \\
Casp11-C254A & No & $<$10 & Minimal \\
Casp11-D285A & No / minimal & $<$15 & Minimal \\
Untransfected + LPS & No & $<$5 & Minimal \\
\bottomrule
\end{tabular}
\end{table}

\subsection{Biochemical evidence for activity-dependent oligomerization}
Size-exclusion chromatography of lysates from LPS-stimulated macrophages revealed that wild-type caspase-11 shifted to a high molecular weight fraction consistent with SMOC formation (Table~\ref{tab:sec}). In the absence of LPS, wild-type caspase-11 remained monomeric. Critically, both C254A and D285A mutants remained in monomeric or low molecular weight fractions even after LPS stimulation, confirming that catalytic activity and autoprocessing are required for oligomeric complex formation.

\begin{table}[h]
\centering
\caption{Oligomeric state of caspase-11 constructs assessed by size-exclusion chromatography.}
\label{tab:sec}
\begin{tabular}{lcc}
\toprule
Construct & LPS treatment & Oligomeric state \\
\midrule
Casp11-WT & + LPS & High MW complex (SMOC) \\
Casp11-WT & -- LPS & Monomeric \\
Casp11-C254A & + LPS & Monomeric / low MW \\
Casp11-D285A & + LPS & Monomeric / low MW \\
\bottomrule
\end{tabular}
\end{table}

\section{Discussion}

Our study provides the first direct visualization of caspase-11 speck formation in macrophages, establishing that the noncanonical inflammasome assembles into a discrete SMOC analogous to the ASC speck of canonical inflammasomes. The key finding is that catalytic activity and intra-molecular autoprocessing are required for this assembly, which challenges the prevailing model of inflammasome formation.

In the canonical inflammasome, sensor proteins (NLRP3, AIM2) oligomerize upon detecting their respective danger signals, recruit the adaptor ASC, and the resulting SMOC activates caspase-1 through proximity-induced dimerization \citep{lu2014unified}. The catalytic activity of caspase-1 is generally thought to be a downstream consequence of assembly, not a requirement for it. Our data suggest that caspase-11 operates differently: its enzymatic activity appears to be required for or coupled to the oligomerization process itself.

We propose a revised model in which LPS binding to the caspase-11 CARD domain initiates a conformational change that enables limited catalytic activity, and the resulting autoprocessing drives a positive feedback loop that promotes further oligomerization into a stable SMOC. This model is supported by the observation that both the catalytic-dead (C254A) and autoprocessing-deficient (D285A) mutants fail to oligomerize, even though they retain the CARD domain required for LPS binding.

These findings have implications for therapeutic targeting of the noncanonical inflammasome. If catalytic activity is required for complex assembly, then catalytic inhibitors would block not only downstream signaling but also the initial assembly step, potentially providing more complete suppression of noncanonical inflammasome activation than approaches targeting the assembled complex.

A limitation of this study is the use of overexpressed, GFP-tagged constructs, which may not fully recapitulate endogenous caspase-11 behavior. Future studies using knock-in tagged alleles at the endogenous locus will be important for validation.

\section{Conclusion}

We demonstrate that caspase-11 forms LPS-induced specks in macrophages that represent noncanonical inflammasome SMOCs. Unexpectedly, both catalytic activity and intra-molecular autoprocessing are required for speck formation and oligomeric complex assembly, revising the standard model of inflammasome formation. These findings establish that caspase-11 is not a passive passenger recruited to a pre-formed complex but actively drives its own inflammasome assembly through enzymatic activity.

\bibliographystyle{plainnat}
\bibliography{references}

\end{document}
